\section{Módulo 3: O Motor de Recomendação (Smart Context) e Compressores Externos}

Com o Módulo 2 encerrado, o RetroForge já sabe responder a uma pergunta fundamental: \textit{``que console é esse arquivo?''}.
O \texttt{ConsoleIdentifier} devolve, com segurança e sem depender de internet, um valor de \texttt{Console} (definido uma única vez em \texttt{models/Console.hpp}, desde o Módulo 0) -- inclusive os dois valores sentinela que existem especificamente para impedir tragédias (US06 e sua extensão): \texttt{Console::DreamcastModificado} e \texttt{Console::PSP\_PBP}.

Mas identificar não é agir. Um \texttt{Console::PS2} sozinho não comprime nada;
alguém precisa traduzir ``isso é um PS2'' em ``rode o \texttt{chdman} com estes argumentos, e não o \texttt{maxcso}''.
Esse ``alguém'' é o Motor de Recomendação -- o \textbf{Smart Context} descrito na Seção 3.2 da documentação v2.0 -- e é o assunto deste módulo.

\begin{CaixaInfo}[title={O que já existe vs. o que este módulo constrói}]
Do Módulo 0 e do Módulo 2, já temos prontos: \texttt{models::Console} (a única fonte de verdade do vocabulário de consoles, reexportada por \texttt{identification::Console} desde o Módulo 0 -- incluindo \texttt{DreamcastModificado}, \texttt{SegaCD}, \texttt{Saturn}, a distinção \texttt{PSP} vs.\ \texttt{PSP\_PBP}, \textbf{e} os dois consoles de cartucho já com detector concreto pronto: \texttt{SNES} (via \texttt{SnesDetector}, um \texttt{IContainerDetector} de extensão \texttt{.smc}/\texttt{.sfc}) e \texttt{Genesis} (via \texttt{GenesisDetector}, um \texttt{IConsoleDetector} de assinatura binária)), \texttt{ConsoleIdentifier}, os detectores de \texttt{PS1/PS2/GameCube/PSP/PBP/Dreamcast/SegaCD/Saturn/SNES/Genesis}, o \texttt{CueSheetResolver}, a \texttt{IdentificationEngineFactory} (Módulo 2), e duas classes de infraestrutura que vamos \textbf{reaproveitar} agora pela primeira vez fora do Módulo 0: \texttt{io::BinaryFileReader} e \texttt{process::ExternalProcessHandle}.

Este módulo é exatamente o que falta escrever e commitar para fechar esse gap: vamos (1) criar a interface \texttt{ICompressionStrategy}, (2) implementar \texttt{ChdmanConverter}, \texttt{DolphinToolConverter} e \texttt{MaxcsoConverter} -- todos por trás de \texttt{ExternalProcessHandle}, (3) implementar \texttt{LibzipPacker}, que \textbf{não} usa subprocesso, (4) construir o \texttt{CompressionRouter} que orquestra tudo isso e bloqueia explicitamente \textbf{os dois} valores sentinela que o Módulo 2 introduziu, e (5) -- diferente da primeira versão deste módulo -- montar essa composição inteira através de uma classe própria, \texttt{CompressionEngineFactory}, em vez de uma função solta.

\begin{CaixaAlerta}[title={Duas correções em relação a uma versão anterior deste módulo}]
Uma versão anterior deste módulo tinha dois problemas que esta revisão corrige:

\begin{enumerate}
    \item \textbf{Bloqueio incompleto:} tratava apenas de \texttt{Console::DreamcastModificado} como caso a bloquear, porque, na época em que foi escrita, o Módulo 2 ainda não havia introduzido \texttt{Console::PSP\_PBP}. O \texttt{CompressionRouter} apresentado a seguir bloqueia \textbf{ambos} os valores sentinela, com o mesmo \texttt{if} explícito e antecipado -- nenhuma rota é registrada para nenhum dos dois, e a ausência de registro nunca é, sozinha, o mecanismo de segurança.
    \item \textbf{Montagem como função solta:} a composição do motor de recomendação vivia em uma função livre, \texttt{construir\_motor\_de\_recomendacao()}, escrita fora de qualquer classe -- exatamente o padrão que o Módulo 2 já havia rejeitado para o Motor de Identificação (Seção sobre \texttt{IdentificationEngineFactory}, que inclusive já previa esta correção: ``o Módulo 3 vai introduzir, no mesmo espírito, uma \texttt{CompressionEngineFactory}''). Esta revisão substitui a função solta por essa classe, na mesma pasta \texttt{bootstrap/} que já hospeda a \texttt{IdentificationEngineFactory}.
\end{enumerate}
\end{CaixaAlerta}

Fora essas duas correções, não precisamos ajustar nenhuma outra divergência de vocabulário: como o Módulo 0 já unificou \texttt{Console} em um único arquivo (\texttt{models/Console.hpp}), \texttt{ConversionJob} e \texttt{ConsoleIdentifier} sempre enxergaram exatamente o mesmo tipo -- inclusive todos os valores introduzidos no Módulo 2.
\end{CaixaInfo}

\subsection{Do ``identificar'' ao ``agir'': por que o roteamento é uma classe separada}

Seria tecnicamente possível colocar um gigantesco \texttt{if/else} dentro do próprio \texttt{main()} do \texttt{cli/}, testando o \texttt{Console} devolvido por \texttt{ConsoleIdentifier} e chamando \texttt{system("chdman ...")} diretamente.
O RetroForge rejeita essa abordagem por três razões, todas já sementadas no Módulo 0:

\begin{itemize}
    \item \textbf{Responsabilidade Única:} \texttt{ConsoleIdentifier} sabe reconhecer um console;
ele não deveria saber nada sobre \texttt{chdman} ou \texttt{maxcso}. Se amanhã o projeto trocar \texttt{maxcso} por outra ferramenta de compressão de PSP, isso não pode exigir tocar em uma linha sequer do Motor de Identificação.
    \item \textbf{Separação Core/Apresentação:} a documentação estabelece que \texttt{cli/} (e a futura \texttt{gui/}, Módulo 11) são camadas de apresentação.
Se a lógica de roteamento vivesse dentro de \texttt{cli/main.cpp}, a futura \texttt{gui/} teria que \textbf{duplicar} essa lógica -- exatamente o tipo de acoplamento que o \texttt{core/} existe para evitar.
    \item \textbf{Testabilidade (Módulo 12):} testar ``um \texttt{Console::PSP} deve acionar o \texttt{MaxcsoConverter}'' -- e ``um \texttt{Console::PSP\_PBP} deve ser bloqueado'' -- é trivial se o roteamento for uma classe isolada do \texttt{core}.
Se essa decisão estivesse espalhada dentro de um \texttt{main()} que também lê argumentos de linha de comando, testá-la exigiria simular todo um processo de CLI.
\end{itemize}

\begin{CaixaAlerta}
Vale repetir, porque é o motivador histórico do projeto inteiro (ver \textit{Motivação e História do Projeto}): o roteamento errado -- mandar um \texttt{.cdi} modificado de Dreamcast para o \texttt{chdman}, ou um \texttt{EBOOT.PBP} já empacotado para o \texttt{maxcso} -- não é um bug cosmético, é o \textbf{mesmo padrão} de bug que \textbf{inspirou} o RetroForge, apenas manifestado em dois consoles diferentes.
O Motor de Recomendação deste módulo é, literalmente, o código que existe para que aquele incidente nunca se repita -- em nenhuma das duas formas que ele pode assumir.
\end{CaixaAlerta}

\subsection{Duas naturezas de compressão, uma única abstração}

Olhando para a Seção 2 da documentação v2.0, as ferramentas de conversão do RetroForge se dividem em dois grupos com naturezas de integração completamente diferentes:

\begin{itemize}
    \item \textbf{\texttt{chdman}, \texttt{maxcso}, \texttt{dolphin-tool}:} são executáveis externos, de terceiros, acionados \textbf{via subprocesso} (a mesma técnica de \texttt{popen}/\texttt{pclose} que o \texttt{ExternalProcessHandle} do Módulo 0 já encapsula).
    \item \textbf{\texttt{libzip}:} é uma biblioteca C, \textbf{compilada e linkada estaticamente} dentro do próprio binário do RetroForge (Seção 2 da documentação v2.0: ``libzip \textit{estático no C++}'').
Não existe processo filho, não existe \texttt{stdout} para ler -- é uma chamada de função direta, dentro do mesmo processo.
\end{itemize}

A primeira tentação de design seria criar duas interfaces -- uma para ``compressores de subprocesso'' e outra para ``compressores de biblioteca''.
Mas isso seria um erro: o \textbf{contrato} que o resto do sistema precisa (``dado um \texttt{ConversionJob}, produza o arquivo comprimido e diga se deu certo'') é \textit{idêntico} nos dois casos.
O que muda é só o \textbf{mecanismo interno} de cada implementação -- exatamente o tipo de detalhe que a Inversão de Dependência existe para esconder de quem consome a abstração.

\begin{itemize}
    \item \textbf{Por que \texttt{ICompressionStrategy} existe:} é o contrato mínimo (Segregação de Interface) que qualquer estratégia de compressão -- baseada em subprocesso ou em biblioteca estática -- deve cumprir.
O \texttt{CompressionRouter} (adiante) depende apenas dela, nunca de \texttt{ChdmanConverter} ou \texttt{LibzipPacker} diretamente (Inversão de Dependência).
    \item \textbf{Como se comunica com outras classes:} é implementada por \texttt{ChdmanConverter}, \texttt{MaxcsoConverter}, \texttt{DolphinToolConverter} (todas usando \texttt{process::ExternalProcessHandle} internamente) e por \texttt{LibzipPacker} (que não usa);
é consumida exclusivamente pelo \texttt{CompressionRouter}.
    \item \textbf{Onde ela mora:} \texttt{core/include/retroforge/compression/ICompressionStrategy.hpp} -- uma pasta nova, irmã de \texttt{identification/}, \texttt{io/}, \texttt{process/} e \texttt{models/}.
\end{itemize}

\begin{CaixaCodigo}
// Arquivo: core/include/retroforge/compression/ICompressionStrategy.hpp
#pragma once

#include "retroforge/models/ConversionJob.hpp"

#include <memory>
#include <string>

namespace retroforge::compression {

// Responsabilidade unica: dado um ConversionJob ja identificado, produzir
// o arquivo comprimido correspondente. NAO sabe como o console foi
// identificado (isso e' problema do Modulo 2) e NAO sabe em que ordem
// os jobs sao executados (isso e' problema do Thread Pool, Modulo 4).
class ICompressionStrategy {
public:
    virtual ~ICompressionStrategy() = default;

    // Executa a compressao de fato. Retorna true em caso de sucesso.
    virtual bool comprimir(const std::shared_ptr<models::ConversionJob>& job) = 0;

    // Extensao do arquivo de SAIDA que esta estrategia produz (".chd",
    // ".cso", ".rvz", ".zip"). Usado pelo CompressionRouter para compor
    // o caminho final e, futuramente, pela Estimativa Preditiva (Modulo 5).
    virtual std::string extensao_saida() const = 0;
};

} // namespace retroforge::compression
\end{CaixaCodigo}

\begin{CaixaAlerta}
Repare que \texttt{<string>} é incluído explicitamente aqui, mesmo que \texttt{models/ConversionJob.hpp} já traga o cabeçalho transitivamente. Depender de uma inclusão transitiva para compilar é uma armadilha comum -- funciona hoje, mas quebra silenciosamente no dia em que alguém reorganizar os includes de \texttt{ConversionJob.hpp}. Cada arquivo deve incluir exatamente o que usa diretamente, nunca menos.
\end{CaixaAlerta}

\begin{CaixaInfo}
Vale uma pausa para apreciar um não-evento: diferente do que aconteceria se o projeto tivesse dois enums \texttt{Console} independentes (um em \texttt{models/}, outro em \texttt{identification/}), \texttt{job->console\_alvo()} já devolve, neste exato momento, o mesmo tipo -- incluindo \texttt{DreamcastModificado}, \texttt{PSP\_PBP}, \texttt{SegaCD}, \texttt{Saturn}, \texttt{SNES} e \texttt{Genesis} -- que \texttt{ConsoleIdentifier::identify()} produziu no Módulo 2. Isso só é verdade porque o Módulo 0 tomou a decisão de extrair \texttt{Console} para \texttt{models/Console.hpp} \textbf{antes} de qualquer código depender dele, e o Módulo 2 amplificou aquele único arquivo, não uma cópia.
O \texttt{CompressionRouter} que vamos construir nesta seção pode, portanto, confiar cegamente em \texttt{job->console\_alvo()} para distinguir rotas específicas, sem nenhuma etapa de conciliação entre vocabulários.
\end{CaixaInfo}

\subsection{As estratégias baseadas em subprocesso: \texttt{ChdmanConverter}, \texttt{DolphinToolConverter} e \texttt{MaxcsoConverter}}

As três ferramentas externas de disco óptico/Nintendo/PSP compartilham o mesmo mecanismo -- abrir um processo filho via \texttt{ExternalProcessHandle} (Módulo 0) e ler sua saída -- mudando apenas o comando montado e a extensão de saída.

\begin{itemize}
    \item \textbf{Por que \texttt{ChdmanConverter} existe:} converte imagens de disco óptico (PS1, PS2, Sega CD, Saturn e Dreamcast limpo) para \texttt{.chd}, delegando o trabalho pesado ao executável \texttt{chdman} via subprocesso.
    \item \textbf{Como se comunica com outras classes:} usa \texttt{process::ExternalProcessHandle} como dependência interna (\texttt{unique\_ptr}, posse exclusiva); recebe um \texttt{std::shared\_ptr<models::ConversionJob>} do \texttt{CompressionRouter} e chama \texttt{set\_progresso} nele conforme lê a saída do processo.
    \item \textbf{Onde ela mora:} \texttt{core/include/retroforge/compression/converters/ChdmanConverter.hpp} e \texttt{core/src/compression/converters/ChdmanConverter.cpp}.
\end{itemize}

\begin{CaixaCodigo}
// Arquivo: core/include/retroforge/compression/converters/ChdmanConverter.hpp
#pragma once

#include "retroforge/compression/ICompressionStrategy.hpp"

namespace retroforge::compression::converters {

class ChdmanConverter : public ICompressionStrategy {
public:
    bool comprimir(const std::shared_ptr<models::ConversionJob>& job) override;
    std::string extensao_saida() const override { return ".chd"; }
};

} // namespace retroforge::compression::converters
\end{CaixaCodigo}

\begin{CaixaCodigo}
// Arquivo: core/src/compression/converters/ChdmanConverter.cpp
#include "retroforge/compression/converters/ChdmanConverter.hpp"

#include "retroforge/process/ExternalProcessHandle.hpp"

#include <exception>
#include <filesystem>
#include <iostream>

namespace retroforge::compression::converters {

namespace fs = std::filesystem;

bool ChdmanConverter::comprimir(const std::shared_ptr<models::ConversionJob>& job) {
    fs::path origem(job->caminho_origem());
    fs::path destino = origem;
    destino.replace_extension(".chd");

    std::string comando =
        "chdman createcd -i \"" + origem.string() + "\" -o \"" + destino.string() + "\"";

    try {
        process::ExternalProcessHandle processo(comando);

        std::string linha;
        while (processo.read_line(linha)) {
            if (job->is_cancelado()) {
                return false;
            }
        }

        job->set_progresso(100);
        return true;

    } catch (const std::exception& erro) {
        std::cerr << "Falha ao converter via chdman: " << erro.what() << "\n";
        return false;
    }
}

} // namespace retroforge::compression::converters
\end{CaixaCodigo}

\texttt{DolphinToolConverter} segue o mesmo molde, com a rota Nintendo habilitando a eliminação de ``junk data'' (Seção 3.2 da documentação v2.0):

\begin{CaixaCodigo}
// Arquivo: core/include/retroforge/compression/converters/DolphinToolConverter.hpp
#pragma once

#include "retroforge/compression/ICompressionStrategy.hpp"

namespace retroforge::compression::converters {

class DolphinToolConverter : public ICompressionStrategy {
public:
    bool comprimir(const std::shared_ptr<models::ConversionJob>& job) override;
    std::string extensao_saida() const override { return ".rvz"; }
};

} // namespace retroforge::compression::converters
\end{CaixaCodigo}

\begin{CaixaCodigo}
// Arquivo: core/src/compression/converters/DolphinToolConverter.cpp
#include "retroforge/compression/converters/DolphinToolConverter.hpp"

#include "retroforge/process/ExternalProcessHandle.hpp"

#include <exception>
#include <filesystem>
#include <iostream>

namespace retroforge::compression::converters {

namespace fs = std::filesystem;

bool DolphinToolConverter::comprimir(const std::shared_ptr<models::ConversionJob>& job) {
    fs::path origem(job->caminho_origem());
    fs::path destino = origem;
    destino.replace_extension(".rvz");

    std::string comando = "dolphin-tool convert -i \"" + origem.string() + "\" -o \"" +
                          destino.string() + "\" -f rvz -b 131072 -c zstd -l 5";

    try {
        process::ExternalProcessHandle processo(comando);

        std::string linha;
        while (processo.read_line(linha)) {
            if (job->is_cancelado()) {
                return false;
            }
        }

        job->set_progresso(100);
        return true;

    } catch (const std::exception& erro) {
        std::cerr << "Falha ao converter via dolphin-tool: " << erro.what() << "\n";
        return false;
    }
}

} // namespace retroforge::compression::converters
\end{CaixaCodigo}

\texttt{MaxcsoConverter} (para \texttt{Console::PSP} -- o \texttt{.iso} cru, \textbf{nunca} \texttt{PSP\_PBP} -- $\rightarrow$ \texttt{.cso}) segue exatamente o mesmo molde:

\begin{CaixaCodigo}
// Arquivo: core/include/retroforge/compression/converters/MaxcsoConverter.hpp
#pragma once

#include "retroforge/compression/ICompressionStrategy.hpp"

namespace retroforge::compression::converters {

class MaxcsoConverter : public ICompressionStrategy {
public:
    bool comprimir(const std::shared_ptr<models::ConversionJob>& job) override;
    std::string extensao_saida() const override { return ".cso"; }
};

} // namespace retroforge::compression::converters
\end{CaixaCodigo}

\begin{CaixaCodigo}
// Arquivo: core/src/compression/converters/MaxcsoConverter.cpp
#include "retroforge/compression/converters/MaxcsoConverter.hpp"

#include "retroforge/process/ExternalProcessHandle.hpp"

#include <exception>
#include <filesystem>
#include <iostream>

namespace retroforge::compression::converters {

namespace fs = std::filesystem;

bool MaxcsoConverter::comprimir(const std::shared_ptr<models::ConversionJob>& job) {
    fs::path origem(job->caminho_origem());
    fs::path destino = origem;
    destino.replace_extension(".cso");

    std::string comando = "maxcso \"" + origem.string() + "\" -o \"" + destino.string() + "\"";

    try {
        process::ExternalProcessHandle processo(comando);

        std::string linha;
        while (processo.read_line(linha)) {
            if (job->is_cancelado()) {
                return false;
            }
        }

        job->set_progresso(100);
        return true;

    } catch (const std::exception& erro) {
        // Mensagem de erro corrigida: esta funcao usa maxcso, nao chdman.
        std::cerr << "Falha ao converter via maxcso: " << erro.what() << "\n";
        return false;
    }
}

} // namespace retroforge::compression::converters
\end{CaixaCodigo}

\subsection{A estratégia em processo: \texttt{LibzipPacker}}

Cartuchos (\texttt{Console::SNES}, \texttt{Console::Genesis} -- Seção 3.2 da documentação v2.0) não passam por nenhum executável externo: o próprio C++ do RetroForge encapsula o arquivo em \texttt{.zip} usando a biblioteca \texttt{libzip}, estaticamente linkada.

\begin{itemize}
    \item \textbf{Por que \texttt{LibzipPacker} existe:} sua única responsabilidade é envolver um arquivo de cartucho em um contêiner \texttt{.zip} válido, usando chamadas diretas da API C do \texttt{libzip} -- sem subprocesso, sem \texttt{popen}, sem depender de nenhum executável externo estar instalado no sistema.
    \item \textbf{Como se comunica com outras classes:} implementa \texttt{ICompressionStrategy}, exatamente como as três classes anteriores -- para o \texttt{CompressionRouter}, ela é indistinguível das demais.
    \item \textbf{Onde ela mora:} \texttt{core/include/retroforge/compression/converters/LibzipPacker.hpp} e \texttt{core/src/compression/converters/LibzipPacker.cpp}.
\end{itemize}

\begin{CaixaCodigo}
// Arquivo: core/include/retroforge/compression/converters/LibzipPacker.hpp
#pragma once

#include "retroforge/compression/ICompressionStrategy.hpp"

namespace retroforge::compression::converters {

class LibzipPacker : public ICompressionStrategy {
public:
    bool comprimir(const std::shared_ptr<models::ConversionJob>& job) override;
    std::string extensao_saida() const override { return ".zip"; }
};

} // namespace retroforge::compression::converters
\end{CaixaCodigo}

\begin{CaixaInfo}[title={Compilando \texttt{libzip} em Windows e Linux}]
O texto original deste módulo deixava a integração com \texttt{libzip} em aberto (\textit{``decisão de build ainda pendente''}). Na prática, a forma mais robusta de resolver isso \textbf{sem} depender de um pacote instalado manualmente pelo usuário em cada sistema operacional é ramificar o \texttt{core/CMakeLists.txt} por plataforma: no Linux, \texttt{pkg-config} localiza o \texttt{libzip-dev} do sistema (instalado via \texttt{apt} no CI, no mesmo espírito do \texttt{clang-format}); no Windows (e macOS), \texttt{find\_package(libzip CONFIG REQUIRED)} espera um pacote gerenciado (tipicamente via \texttt{vcpkg}), que já expõe um alvo \texttt{CMake} nativo. Em ambos os casos, o projeto termina com um único alvo lógico, \texttt{libzip::zip}, que o restante do \texttt{CMakeLists.txt} linka sem precisar saber qual caminho foi usado para resolvê-lo:

\begin{CaixaCodigo}
# Arquivo: core/CMakeLists.txt (trecho adicionado no topo)
if(UNIX AND NOT APPLE)
    find_package(PkgConfig REQUIRED)
    pkg_check_modules(LIBZIP REQUIRED IMPORTED_TARGET libzip)

    if(TARGET PkgConfig::LIBZIP AND NOT TARGET libzip::zip)
        add_library(libzip::zip ALIAS PkgConfig::LIBZIP)
    endif()
else()
    # Windows e macOS usam o find_package nativo / vcpkg
    find_package(libzip CONFIG REQUIRED)
endif()

# ... (targets existentes) ...

target_link_libraries(retroforge_core PRIVATE
    libzip::zip
)
\end{CaixaCodigo}

No \texttt{.github/workflows/cmake-single-platform.yml}, o passo de instalação de dependências do sistema precisa crescer para incluir \texttt{pkg-config} e \texttt{libzip-dev} (Ubuntu, runner atual do CI), antes mesmo da checagem de formatação:

\begin{CaixaCodigo}
# Arquivo: .github/workflows/cmake-single-platform.yml (trecho revisado)
- name: Install System Dependencies
  run: |
    sudo apt-get update
    sudo apt-get install -y clang-format pkg-config libzip-dev

- name: Check Code Formatting
  run: |
    find core cli -type f \( -name "*.cpp" -o -name "*.hpp" -o -name "*.h" \) | xargs -r clang-format --dry-run --Werror
\end{CaixaCodigo}

Isso antecipa, de forma concreta, o multiplataforma que o Módulo 13 (CI/CD Avançado) só formaliza mais tarde -- mas resolver a dependência de \texttt{libzip} sem travar o Windows já é necessário agora, e não faz sentido adiar.
\end{CaixaInfo}

\begin{CaixaCodigo}
// Arquivo: core/src/compression/converters/LibzipPacker.cpp
#include "retroforge/compression/converters/LibzipPacker.hpp"

#include <filesystem>
#include <iostream>
#include <memory>

#include <zip.h>

namespace retroforge::compression::converters {

namespace fs = std::filesystem;

namespace {

// zip_discard() (nao zip_close()) e' o deleter correto para o caminho de
// RAII "normal": zip_close() so' e' chamado explicitamente apos o sucesso
// de zip_file_add, e so' ENTAO o handle e' liberado do unique_ptr via
// release() -- ver comprimir() abaixo. Se qualquer 'return false' anterior
// disparar, o destrutor chama zip_discard(), que fecha o arquivo SEM
// gravar as mudancas -- nunca deixando um .zip corrompido/parcial no disco.
struct ZipDiscardDeleter {
    void operator()(zip_t* z) const {
        if (z) {
            zip_discard(z);
        }
    }
};
using ScopedZipHandle = std::unique_ptr<zip_t, ZipDiscardDeleter>;

} // namespace

bool LibzipPacker::comprimir(const std::shared_ptr<models::ConversionJob>& job) {
    fs::path origem(job->caminho_origem());
    fs::path destino = origem;
    destino.replace_extension(".zip");

    int codigo_erro = 0;

#if defined(_WIN32)
    // No Windows, std::filesystem::path armazena wchar_t internamente; a
    // API C do libzip espera UTF-8 em todas as plataformas -- u8string()
    // faz essa conversao explicitamente, em vez de confiar em .string(),
    // que no Windows usaria a codepage local e corromperia nomes com
    // acentos ou caracteres fora do ASCII.
    std::string str_destino = reinterpret_cast<const char*>(destino.u8string().c_str());
    std::string str_origem = reinterpret_cast<const char*>(origem.u8string().c_str());
    std::string str_nome_arquivo =
        reinterpret_cast<const char*>(origem.filename().u8string().c_str());
#else
    std::string str_destino = destino.string();
    std::string str_origem = origem.string();
    std::string str_nome_arquivo = origem.filename().string();
#endif

    ScopedZipHandle zip_guard(
        zip_open(str_destino.c_str(), ZIP_CREATE | ZIP_TRUNCATE, &codigo_erro));

    if (!zip_guard) {
        std::cerr << "Falha ao criar .zip (codigo libzip: " << codigo_erro << ")\n";
        return false;
    }

    auto* fonte = zip_source_file(zip_guard.get(), str_origem.c_str(), 0, 0);

    if (!fonte) {
        std::cerr << "Falha ao ler arquivo de origem para empacotamento.\n";
        return false;
    }

    if (zip_file_add(zip_guard.get(), str_nome_arquivo.c_str(), fonte, ZIP_FL_OVERWRITE) < 0) {
        zip_source_free(fonte);
        std::cerr << "Falha ao adicionar arquivo dentro do .zip.\n";
        return false;
    }

    // release() transfere a posse manualmente para fora do unique_ptr:
    // a partir daqui, SOMOS NOS quem decide como fechar o handle -- por
    // zip_close() (grava de verdade), nunca mais por zip_discard().
    zip_t* zip_raw = zip_guard.release();

    if (zip_close(zip_raw) < 0) {
        std::cerr << "Falha ao gravar arquivo .zip no disco.\n";
        return false;
    }

    job->set_progresso(100);
    return true;
}

} // namespace retroforge::compression::converters
\end{CaixaCodigo}

\subsection{O coração do módulo: \texttt{CompressionRouter}}

Chegamos à classe que dá nome ao módulo.
O \texttt{CompressionRouter} é o Smart Context da documentação: ele recebe um \texttt{ConversionJob} já identificado e decide, com base no \texttt{Console} alvo, qual \texttt{ICompressionStrategy} executar -- ou se a execução deve ser \textbf{bloqueada}. A classe em si não muda em relação à primeira versão deste módulo -- o que muda, na próxima seção, é \textbf{quem a monta}.

\begin{itemize}
    \item \textbf{Por que \texttt{CompressionRouter} existe:} é a única classe que conhece a tabela de roteamento completa.
Nenhuma outra classe do projeto -- nem os conversores, nem o identificador -- precisa saber essa tabela inteira.
    \item \textbf{Como se comunica com outras classes:} depende apenas da abstração \texttt{ICompressionStrategy} (Inversão de Dependência), recebe estratégias via \texttt{register\_strategy} (Princípio Aberto/Fechado) e é consumido pela camada de apresentação depois que \texttt{ConsoleIdentifier::identify} já devolveu um \texttt{Console}.
    \item \textbf{Onde ela mora:} \texttt{core/include/retroforge/compression/CompressionRouter.hpp} e \texttt{core/src/compression/CompressionRouter.cpp} -- direto dentro de \texttt{compression/}, ao lado de \texttt{ICompressionStrategy.hpp} (e não dentro de \texttt{converters/}, porque ela \textit{orquestra} estratégias, não é uma).
\end{itemize}

\begin{CaixaCodigo}
// Arquivo: core/include/retroforge/compression/CompressionRouter.hpp
#pragma once

#include <memory>
#include <unordered_map>

#include "retroforge/compression/ICompressionStrategy.hpp"
#include "retroforge/models/Console.hpp"
#include "retroforge/models/ConversionJob.hpp"

namespace retroforge::compression {

enum class ResultadoRoteamento {
    ConversaoIniciada,
    BloqueadoPorSeguranca,
    ConsoleNaoSuportado
};

class CompressionRouter {
public:
    void register_strategy(models::Console console,
                            std::shared_ptr<ICompressionStrategy> estrategia);

    ResultadoRoteamento rotear(const std::shared_ptr<models::ConversionJob>& job);

private:
    std::unordered_map<models::Console, std::shared_ptr<ICompressionStrategy>>
        estrategias_;
};

} // namespace retroforge::compression
\end{CaixaCodigo}

\begin{CaixaCodigo}
// Arquivo: core/src/compression/CompressionRouter.cpp
#include "retroforge/compression/CompressionRouter.hpp"

namespace retroforge::compression {

using models::Console;

void CompressionRouter::register_strategy(Console console,
                                           std::shared_ptr<ICompressionStrategy> estrategia) {
    estrategias_[console] = std::move(estrategia);
}

ResultadoRoteamento CompressionRouter::rotear(const std::shared_ptr<models::ConversionJob>& job) {
    // A trava de seguranca vem ANTES de qualquer consulta ao mapa de
    // estrategias -- DreamcastModificado (.cdi) e PSP_PBP (EBOOT.PBP ja
    // empacotado) nunca sao alcancados, mesmo que, por engano, alguem
    // registre uma estrategia para eles no futuro.
    if (job->console_alvo() == Console::DreamcastModificado ||
        job->console_alvo() == Console::PSP_PBP) {
        return ResultadoRoteamento::BloqueadoPorSeguranca;
    }

    auto encontrado = estrategias_.find(job->console_alvo());
    if (encontrado == estrategias_.end()) {
        return ResultadoRoteamento::ConsoleNaoSuportado;
    }

    bool sucesso = encontrado->second->comprimir(job);
    return sucesso ? ResultadoRoteamento::ConversaoIniciada
                    : ResultadoRoteamento::ConsoleNaoSuportado;
}

} // namespace retroforge::compression
\end{CaixaCodigo}

\begin{CaixaAlerta}
Repare que o bloqueio de \texttt{DreamcastModificado} \textbf{e} de \texttt{PSP\_PBP} é feito com um único \texttt{if} \textbf{explícito}, checado antes de qualquer coisa -- não com um simples ``esquecer de registrar uma estratégia para esses valores''. A trava explícita, verificada antes do \texttt{unordered\_map} e cobrindo \textbf{ambos} os valores sentinela, é impossível de contornar por acidente -- mesmo que o \texttt{enum class Console} ganhe um terceiro valor sentinela em módulos futuros, o padrão a seguir já está estabelecido aqui.
\end{CaixaAlerta}

\subsection{Montando o motor completo: \texttt{CompressionEngineFactory}}
\label{sec:compression-engine-factory}

Assim como o Motor de Identificação (Módulo 2) parou de depender de uma função solta assim que \texttt{IdentificationEngineFactory} nasceu, o Motor de Recomendação merece exatamente o mesmo tratamento. Decidir quais estratégias existem, para quais consoles cada uma é registrada, e em que ordem isso acontece é, tanto quanto a montagem do Motor de Identificação, uma decisão de regra de negócio -- não um detalhe de demonstração que deveria viver dentro de um \texttt{main()} ou de uma função livre em algum arquivo de cabeçalho solto.

\begin{itemize}
    \item \textbf{Por que \texttt{CompressionEngineFactory} existe:} sua única responsabilidade é saber \textit{compor} um \texttt{CompressionRouter} completo -- registrar cada \texttt{ICompressionStrategy} concreta para os consoles corretos. Ela não sabe \textbf{como} nenhuma estratégia comprime -- isso continua sendo, exclusivamente, problema de cada \texttt{ICompressionStrategy} concreta.
    \item \textbf{Como se comunica com outras classes:} depende de \texttt{CompressionRouter} e de todos os conversores concretos de \texttt{compression::converters} (ela é, de propósito, a única classe do projeto que precisa conhecer essa lista inteira). É consumida pela camada de apresentação (\texttt{cli/}, Módulo 11, e futuramente \texttt{gui/}) e pelo \texttt{BatchConversionManager} (Módulo 4), que passam a chamar apenas \texttt{CompressionEngineFactory::criar()} em vez de repetir a lista de \texttt{register\_strategy} em cada ponto de entrada do programa.
    \item \textbf{Onde ela mora:} \texttt{core/include/retroforge/bootstrap/CompressionEngineFactory.hpp} e \texttt{core/src/bootstrap/CompressionEngineFactory.cpp} -- a mesma pasta \texttt{bootstrap/} que o Módulo 2 já criou para a \texttt{IdentificationEngineFactory}, exatamente como aquele módulo já previa (\textit{``o Módulo 3 vai introduzir, no mesmo espírito, uma \texttt{CompressionEngineFactory}, sem que esta classe [\texttt{IdentificationEngineFactory}] precise ser tocada''}) -- e, de fato, nenhuma linha de \texttt{IdentificationEngineFactory.hpp/.cpp} muda por causa deste módulo.
\end{itemize}

\begin{CaixaCodigo}
// Arquivo: core/include/retroforge/bootstrap/CompressionEngineFactory.hpp
#pragma once

#include "retroforge/compression/CompressionRouter.hpp"

namespace retroforge::bootstrap {

// Responsabilidade unica: saber MONTAR o Motor de Recomendacao completo --
// registrar, para cada Console suportado, a ICompressionStrategy correta.
// NAO decide COMO uma estrategia comprime -- isso continua sendo
// responsabilidade exclusiva de cada conversor concreto. Espelha
// IdentificationEngineFactory (Modulo 2) tanto na intencao quanto na
// pasta onde mora -- bootstrap/ e' reservada a classes de COMPOSICAO,
// nunca a regra de negocio em si.
class CompressionEngineFactory {
public:
    // Constroi um CompressionRouter pronto para uso, com todas as rotas
    // seguras ja registradas. DreamcastModificado e PSP_PBP nunca sao
    // registrados aqui -- e nao precisam ser, ja que CompressionRouter::
    // rotear() os bloqueia antes de qualquer consulta ao mapa interno.
    static compression::CompressionRouter criar();
};

} // namespace retroforge::bootstrap
\end{CaixaCodigo}

\begin{CaixaCodigo}
// Arquivo: core/src/bootstrap/CompressionEngineFactory.cpp
#include "retroforge/bootstrap/CompressionEngineFactory.hpp"

#include "retroforge/compression/converters/ChdmanConverter.hpp"
#include "retroforge/compression/converters/DolphinToolConverter.hpp"
#include "retroforge/compression/converters/LibzipPacker.hpp"
#include "retroforge/compression/converters/MaxcsoConverter.hpp"

#include <memory>

namespace retroforge::bootstrap {

compression::CompressionRouter CompressionEngineFactory::criar() {
    compression::CompressionRouter roteador;

    using models::Console;

    // Discos opticos -> chdman. O mesmo shared_ptr e' reaproveitado para
    // as cinco rotas: ChdmanConverter nao guarda estado proprio entre
    // chamadas, entao compartilhar a instancia evita alocacoes redundantes.
    auto chdman = std::make_shared<compression::converters::ChdmanConverter>();
    roteador.register_strategy(Console::PS1, chdman);
    roteador.register_strategy(Console::PS2, chdman);
    roteador.register_strategy(Console::SegaCD, chdman);
    roteador.register_strategy(Console::Saturn, chdman);
    roteador.register_strategy(Console::Dreamcast, chdman); // .gdi limpo

    // GameCube -> dolphin-tool.
    roteador.register_strategy(
        Console::GameCube, std::make_shared<compression::converters::DolphinToolConverter>());

    // PSP cru -> maxcso. Console::PSP_PBP NUNCA e' registrado aqui,
    // propositalmente (ver CompressionRouter::rotear).
    roteador.register_strategy(Console::PSP,
                                std::make_shared<compression::converters::MaxcsoConverter>());

    // Cartuchos -> libzip. Diferente da primeira versao deste modulo
    // (quando LibzipPacker nao tinha nenhum detector correspondente
    // ainda), o Modulo 2 ja entrega SnesDetector (Console::SNES, via
    // extensao .smc/.sfc) e GenesisDetector (Console::Genesis, via
    // assinatura binaria) prontos -- a lacuna documentada la esta
    // fechada aqui, com o mesmo shared_ptr compartilhado entre as duas
    // rotas, pelo mesmo motivo do chdman acima.
    auto libzip = std::make_shared<compression::converters::LibzipPacker>();
    roteador.register_strategy(Console::SNES, libzip);
    roteador.register_strategy(Console::Genesis, libzip);

    // Nem Console::DreamcastModificado nem Console::PSP_PBP sao
    // registrados aqui -- a trava explicita dentro de rotear() ja cobre
    // os dois, e essa ausencia de registro e' redundante de proposito.
    return roteador;
}

} // namespace retroforge::bootstrap
\end{CaixaCodigo}

Com a fábrica pronta, qualquer ponto de entrada -- o \texttt{main()} de demonstração deste módulo, o futuro \texttt{cli/main.cpp} (Módulo 11), ou o \texttt{BatchConversionManager} (Módulo 4) -- monta o motor completo com uma única chamada:

\begin{CaixaCodigo}
#include <iostream>

#include "retroforge/bootstrap/CompressionEngineFactory.hpp"
#include "retroforge/bootstrap/IdentificationEngineFactory.hpp"

using namespace retroforge;

int main() {
    auto identificador = bootstrap::IdentificationEngineFactory::criar(); // Modulo 2
    auto roteador = bootstrap::CompressionEngineFactory::criar();         // Modulo 3

    auto console = identificador.identify("jogo.iso");

    auto job = std::make_shared<models::ConversionJob>("jogo.iso", console);
    auto resultado = roteador.rotear(job);

    switch (resultado) {
        case compression::ResultadoRoteamento::ConversaoIniciada:
            std::cout << "Conversao concluida com sucesso.\n";
            break;
        case compression::ResultadoRoteamento::BloqueadoPorSeguranca:
            std::cout << "Bloqueado por seguranca (Dreamcast .cdi ou PSP EBOOT.PBP).\n";
            break;
        case compression::ResultadoRoteamento::ConsoleNaoSuportado:
            std::cout << "Nenhuma rota registrada para este console, ou falha na compressao.\n";
            break;
    }

    return 0;
}
\end{CaixaCodigo}

\begin{CaixaInfo}
Note a simetria completa com o Módulo 2: assim como \texttt{cli/main.cpp} nunca vai precisar conhecer a lista inteira de detectores (isso é responsabilidade exclusiva de \texttt{IdentificationEngineFactory}), ele também nunca vai precisar conhecer a lista inteira de conversores (agora responsabilidade exclusiva de \texttt{CompressionEngineFactory}). As duas fábricas moram lado a lado em \texttt{bootstrap/}, cada uma cuidando de compor exatamente um motor -- nenhuma fábrica genérica que tentasse montar os dois de uma vez, o que a forçaria a depender de \texttt{identification/} \textbf{e} \texttt{compression/} simultaneamente, violando a Segregação de Interface no nível de módulo (o mesmo princípio já aplicado no Módulo 2 para justificar por que \texttt{bootstrap/} não ganha uma única fábrica ``faz-tudo'').
\end{CaixaInfo}

\subsection{\texttt{LibzipPacker} não fica mais sem rota}

A primeira versão deste módulo dedicava uma seção inteira a explicar, honestamente, por que \texttt{LibzipPacker} não tinha nenhuma rota registrada: o Módulo 2, na época, ainda não tinha nenhum \texttt{Console::SNES} ou \texttt{Console::MegaDrive} com detector concreto por trás. Isso mudou -- o Módulo 2 já formalizou \texttt{SnesDetector} (contêiner, extensão \texttt{.smc}/\texttt{.sfc}, produzindo \texttt{Console::SNES}) e \texttt{GenesisDetector} (assinatura binária, produzindo \texttt{Console::Genesis}) -- e a \texttt{CompressionEngineFactory} desta seção já fecha essa lacuna, registrando \texttt{LibzipPacker} para os dois.

\begin{CaixaAlerta}
Vale registrar a diferença de nomenclatura entre o vocabulário do \texttt{enum class Console} (que usa \texttt{Genesis}, o nome ocidental do console) e a documentação original do projeto (que menciona ``\texttt{MegaDrive}'' em alguns trechos -- o nome usado fora da América do Norte para o mesmo hardware). Isso não é uma divergência de dois tipos diferentes -- é o mesmo console, um único valor de enum, com um nome de identificador escolhido e usado de forma consistente em todo o código (\texttt{GenesisDetector}, \texttt{Console::Genesis}). Documentação e código descrevendo a mesma coisa com sinônimos diferentes é aceitável; \textbf{dois valores de enum} para a mesma plataforma não seria -- e este projeto tem só um.
\end{CaixaAlerta}

\subsection{Arquitetura de Pastas: o que este módulo acrescenta}

O Módulo 0 já havia reservado o terreno (``cada subpasta corresponde a uma área de responsabilidade''). Este módulo acrescenta \texttt{compression/} (a primeira pasta inteiramente nova de regra de negócio desde \texttt{identification/}) e estende \texttt{bootstrap/}, que o Módulo 2 já havia criado:

\begin{CaixaArvore}
core/
|-- include/retroforge/
|   |-- io/                              [Modulo 0, sem alteracoes]
|   |-- process/                         [Modulo 0, reaproveitado neste modulo]
|   |-- identification/                  [Modulo 0 e 2, sem alteracoes neste modulo]
|   |-- models/                          [Modulo 0/2, sem alteracoes]
|   |-- bootstrap/
|   |   |-- IdentificationEngineFactory.hpp   [Modulo 2, sem alteracoes]
|   |   `-- CompressionEngineFactory.hpp      [NOVO neste modulo]
|   `-- compression/                     [NOVO neste modulo]
|       |-- ICompressionStrategy.hpp     [NOVO]
|       |-- CompressionRouter.hpp        [NOVO]
|       `-- converters/                  [NOVO -- espelha identification/detectors/]
|           |-- ChdmanConverter.hpp      [NOVO]
|           |-- DolphinToolConverter.hpp [NOVO]
|           |-- MaxcsoConverter.hpp      [NOVO]
|           `-- LibzipPacker.hpp         [NOVO]
`-- src/
    |-- bootstrap/
    |   |-- IdentificationEngineFactory.cpp   [Modulo 2, sem alteracoes]
    |   `-- CompressionEngineFactory.cpp      [NOVO neste modulo]
    `-- compression/                     [NOVO neste modulo]
        |-- CompressionRouter.cpp        [NOVO]
        `-- converters/                  [NOVO -- NAO "detectors/", ver alerta acima]
            |-- ChdmanConverter.cpp      [NOVO]
            |-- DolphinToolConverter.cpp [NOVO]
            |-- MaxcsoConverter.cpp      [NOVO]
            `-- LibzipPacker.cpp         [NOVO]
\end{CaixaArvore}

\begin{CaixaInfo}
Graças ao \texttt{file(GLOB\_RECURSE ... CONFIGURE\_DEPENDS)} já configurado no Módulo 1 para \texttt{core/CMakeLists.txt}, toda a pasta \texttt{compression/} entra automaticamente no build de \texttt{retroforge\_core} assim que os arquivos forem salvos em disco -- nenhuma linha do \texttt{CMakeLists.txt} raiz ou de \texttt{core/CMakeLists.txt} precisa ser tocada por causa da nova subpasta. A única linha de build genuinamente pendente é a que resolve a dependência de \texttt{libzip} (Seção sobre \texttt{LibzipPacker} acima) -- essa, sim, exige uma edição real em \texttt{core/CMakeLists.txt} e no workflow de CI, porque \texttt{libzip} não é header-only como o \texttt{CLI11}.
\end{CaixaInfo}

\subsection{Exercícios Práticos do Módulo 3}

\begin{CaixaDesafio}
Lembrete das etiquetas: \textbf{[projeto]} = vira arquivo/pasta permanente do repositório;
\textbf{[laboratório]} = validação prática descartável (não precisa sobreviver no repositório depois de feito).

\begin{enumerate}[label=\textbf{3.\arabic*}]
    \item \textbf{[projeto]} Implemente \texttt{CompressionEngineFactory.hpp}/\texttt{.cpp} conforme apresentado nesta seção, dentro de \texttt{bootstrap/}, ao lado de \texttt{IdentificationEngineFactory}. Se você já tinha uma função solta \texttt{construir\_motor\_de\_recomendacao()} em algum arquivo de demonstração, remova-a -- a fábrica é o único ponto de composição do motor de recomendação a partir de agora.
    \item \textbf{[laboratório]} Sem rodar um jogo real: crie um \texttt{main()} temporário que monte o \texttt{CompressionRouter} via \texttt{CompressionEngineFactory::criar()}, construa manualmente seis \texttt{ConversionJob} fictícios (\texttt{Console::PS2}, \texttt{Console::GameCube}, \texttt{Console::SNES}, \texttt{Console::Genesis}, \texttt{Console::DreamcastModificado} e \texttt{Console::PSP\_PBP}) e confirme, olhando apenas o valor de \texttt{ResultadoRoteamento} devolvido, que os dois últimos são bloqueados \textbf{antes} de qualquer tentativa de chamar um subprocesso ou tocar o \texttt{libzip}. Este \texttt{main()} é descartável -- será formalizado como teste automatizado de verdade no Módulo 12.
    \item \textbf{[projeto]} Confirme que \texttt{core/CMakeLists.txt} resolve \texttt{libzip} corretamente tanto no seu ambiente Linux (via \texttt{pkg-config}) quanto -- se você tiver acesso a uma máquina ou VM Windows -- via \texttt{vcpkg}/\texttt{find\_package(libzip CONFIG REQUIRED)}. Atualize também o workflow de CI para instalar \texttt{pkg-config} e \texttt{libzip-dev} antes da checagem de formatação.
    \item \textbf{[laboratório]} Escreva um teste manual que chame \texttt{LibzipPacker::comprimir} diretamente sobre um \texttt{ConversionJob} fictício apontando para um arquivo de texto qualquer, e confirme que um \texttt{.zip} válido é gerado ao final. Em seguida, force uma falha deliberada (por exemplo, apontando para um arquivo que não existe) e confirme que \textbf{nenhum} arquivo \texttt{.zip} corrompido ou parcial fica no disco -- prova de que o \texttt{zip\_discard()} no destrutor de \texttt{ScopedZipHandle} está funcionando como esperado.
    \item \textbf{[laboratório]} Corrija, no seu clone, a mensagem de erro do \texttt{catch} em \texttt{MaxcsoConverter::comprimir} (Seção sobre os conversores de subprocesso) caso ela ainda mencione ``chdman'' em vez de ``maxcso''. Depois, revise \texttt{ChdmanConverter} e \texttt{DolphinToolConverter} à procura do mesmo tipo de mensagem copiada por engano.
    \item \textbf{[laboratório]} Confirme que os arquivos \texttt{.cpp} deste módulo estão em \texttt{core/src/compression/converters/}, e não em \texttt{core/src/compression/detectors/}. Se encontrar a pasta errada no seu clone, mova os arquivos e confirme que o projeto ainda compila (o \texttt{CONFIGURE\_DEPENDS} do Módulo 1 vai encontrá-los automaticamente no novo caminho).
    \item \textbf{[laboratório]} Compare, com um pequeno teste manual, o comportamento de \texttt{CompressionRouter::rotear} quando um \texttt{ConversionJob} é criado com \texttt{Console::Desconhecido}. Confirme que o resultado é \texttt{ResultadoRoteamento::ConsoleNaoSuportado}, e não um bloqueio de segurança -- três situações semanticamente diferentes, e vale confirmar que o código realmente as distingue, inclusive entre os dois casos de bloqueio (\texttt{DreamcastModificado} e \texttt{PSP\_PBP}).
\end{enumerate}
\end{CaixaDesafio}

\begin{CaixaResumo}
Neste módulo você aprendeu:
\begin{itemize}
    \item Por que o roteamento de compressão (Smart Context) merece ser uma classe própria, separada tanto da identificação (Módulo 2) quanto da apresentação (\texttt{cli/}/\texttt{gui/}).
    \item Como uma única interface (\texttt{ICompressionStrategy}) pode, corretamente, unificar duas naturezas de integração completamente diferentes -- subprocesso e biblioteca estática -- porque o contrato observável é idêntico, mesmo com mecanismos internos radicalmente distintos.
    \item Como reaproveitar, sem alterar, duas classes do Módulo 0 (\texttt{process::ExternalProcessHandle} e o \texttt{shared\_ptr<ConversionJob>}) em um contexto totalmente novo.
    \item Que \texttt{Console::DreamcastModificado} e \texttt{Console::PSP\_PBP} são, arquiteturalmente, o mesmo tipo de problema -- e por que a trava explícita, verificada antes do \texttt{unordered\_map}, cobre ambos.
    \item Por que a \textbf{montagem} do motor de recomendação merece sua própria classe -- \texttt{CompressionEngineFactory}, na pasta \texttt{bootstrap/} já criada pelo Módulo 2 -- em vez de viver como função solta: mesma motivação, mesmo padrão, e mesma pasta que \texttt{IdentificationEngineFactory} já estabeleceu, evitando duplicação futura entre \texttt{cli/} e \texttt{gui/} e mantendo a composição testável isoladamente (Módulo 12).
    \item Que \texttt{LibzipPacker} não fica mais sem rota registrada: os detectores de cartucho que o Módulo 2 já entrega (\texttt{SnesDetector}, \texttt{GenesisDetector}) fecham a lacuna que a primeira versão deste módulo apenas documentava como pendente.
    \item Como resolver \texttt{libzip} de forma multiplataforma desde já (\texttt{pkg-config} no Linux, \texttt{find\_package}/\texttt{vcpkg} no Windows), unificados atrás de um único alvo lógico \texttt{libzip::zip} -- e por que o padrão RAII de \texttt{LibzipPacker} usa \texttt{zip\_discard()} como destrutor padrão, promovendo para \texttt{zip\_close()} explícito só no caminho de sucesso confirmado.
\end{itemize}

Com o console identificado \textbf{e} roteado com segurança para a ferramenta certa (ou explicitamente bloqueado, nos dois casos sentinela) -- agora incluindo cartuchos de verdade --, falta resolver um problema que a documentação (Seção 3.4) já anunciou desde o início: converter 50 jogos, um de cada vez, de forma sequencial, é lento e desperdiça os múltiplos núcleos de qualquer máquina moderna. O Módulo 4 vai introduzir o Thread Pool customizado do RetroForge -- e é exatamente o \texttt{CompressionRouter} deste módulo, montado agora por \texttt{CompressionEngineFactory}, que cada \textit{worker} da piscina de threads vai chamar, em paralelo, para múltiplos \texttt{ConversionJob} simultâneos.
\end{CaixaResumo}
